\subsection{EXAMPLES OF SEMI-SIMPLE LIE ALGEBRAS}

PROPOSITION 8. Let V be a finite-dimensional vector space. Then \(\mathfrak{gl}(V)\) is reductive. its centre is the set of homotheties of V, its derived algebra is \(\mathfrak{sl}(V)\) and the latter is semi-simple.

The identity representation of \(\mathfrak{gl}(\mathsf{V})\) is simple, hence \(\mathfrak{gl}(\mathsf{V})\) is reductive and therefore \(\mathfrak{gl}(\mathsf{V})\) is the direct sum of its centre \(c\) and its derived algebra \(\mathcal{D}(\mathfrak{gl}(\mathsf{V}))\) . The centre \(c\) is the set of homotheties (Algebra, Chapter II, §2, no. 5, Corollary 1 to Proposition 5). Clearly \(\mathcal{D}(\mathfrak{gl}(\mathsf{V})) \subset \mathfrak{sl}(\mathsf{V})\) . As \(\mathfrak{sl}(\mathsf{V}) \cap c = \{0\}\) , \(\mathcal{D}(\mathfrak{gl}(\mathsf{V})) = \mathfrak{sl}(\mathsf{V})\) . Hence \(\mathfrak{sl}(\mathsf{V})\) is semi-simple.

Example. We identify \(\mathfrak{sl}(\mathbf{K}^{2})\) with the Lie algebra of matrices of order 2 and trace zero. We write

\[
X = \left( \begin{array}{c c} 0 & 1 \\ 0 & 0 \end{array} \right) \qquad Y = \left( \begin{array}{c c} 0 & 0 \\ 1 & 0 \end{array} \right) \qquad H = \left( \begin{array}{c c} 1 & 0 \\ 0 & - 1 \end{array} \right).
\]

Then \(X, Y, H\) form a basis of \(\mathfrak{sl}(\mathbf{K}^2)\) and

\[
[ H, X ] = 2 X \quad [ H, Y ] = - 2 Y \quad [ X, Y ] = H.
\]

As an algebra of dimension 1 or 2 is not semi-simple (no. 1, Remark 1), \(\mathfrak{sl}(\mathbf{K}^2)\) is simple. In fact, \(\mathfrak{sl}(\mathbf{V})\) is simple for \(\dim V \geqslant 2\) , as we shall see later (cf.~also Exercises 21 and 24).

PROPOSITION 9. Let \(\mathrm{V}\) be a vector space of finite dimension \(n\) over \(\mathbf{K}\) and \(\beta\) a nondegenerate symmetric (resp. alternating) bilinear form on \(\mathrm{V}\) . Let \(\mathfrak{g}\) be the Lie algebra consisting of the \(x \in \mathfrak{gl}(\mathrm{V})\) such that \(\beta (xm, m') + \beta (m, xm') = 0\) for all \(m, m'\) in \(\mathrm{V}\) . Then \(\mathfrak{g}\) is reductive; \(\mathfrak{g}\) is even semi-simple except in the case where \(\beta\) is symmetric and \(n = 2\) .

For all \(u \in \mathfrak{gl}(\mathrm{V})\) let \(u^*\) denote its adjoint relative to \(\beta\) ; then \(\operatorname{Tr}(u) = \operatorname{Tr}(u^*)\) by Proposition 7 of Algebra, Chapter IX, § 1, no. 8. The condition

\[
\beta (u m, m ^ {\prime}) + \beta (m, u m ^ {\prime}) = 0
\]

for all \(m\) , \(m'\) in V means that \(u + u^{*} = 0\) . In particular, if \(v \in \mathfrak{gl}(V)\) then \((v - v^{*})^{*} = v^{*} - v\) and hence \(v - v^{*} \in \mathfrak{g}\) . Then let \(u\) be an element of \(\mathfrak{g}\) orthogonal to \(\mathfrak{g}\) with respect to the bilinear form \(\phi\) associated with the identity representation of \(\mathfrak{g}\) . For all \(v \in \mathfrak{gl}(V)\) , \(\operatorname{Tr} u(v - v^{*}) = 0\) , hence

\[
\operatorname{Tr} (u v) = \operatorname{Tr} (u v ^ {*}) = \operatorname{Tr} (u v ^ {*}) ^ {*} = \operatorname{Tr} (v u ^ {*}) = - \operatorname{Tr} (v u) = - \operatorname{Tr} (u v)
\]

and hence \(\operatorname{Tr}(u v) = 0\) . It follows that \(u = 0\) , so that \(\phi\) is non-degenerate. Hence \(\mathfrak{g}\) is reductive (Proposition 5). It remains to show that the centre of \(\mathfrak{g}\) is zero (except when \(\beta\) is symmetric and \(n = 2\) ). By extending the base field, we can assume that \(\mathbf{K}\) is algebraically closed.

\begin{enumerate}
\def\labelenumi{(\alph{enumi})}
\item
  When \(\beta\) is symmetric, it can be identified with the bilinear form on \(\mathbf{K}^n\) with matrix \(I_{\mathrm{n}}\) with respect to the canonical basis (Algebra, Chapter IX, §6, Corollary 1 to Theorem 1). Under these conditions \(\mathfrak{g}\) is identified with the Lie algebra of skew-symmetric matrices (§3, no. 4, Example 1). Let \(U = (u_{ij}) \in \mathfrak{g}\) ; we use the fact that \(U\) commutes with the matrix \((v_{ij}) \in \mathfrak{g}\) all of whose elements are zero except \(v_{i_0j_0}\) and \(v_{j_0i_0}\) ( \(i_0 \neq j_0\) ) which are equal respectively to 1 and -1. We find that \(u_{i_0j} = u_{j_0j} = u_{ii_0} = u_{ij_0} = 0\) for \(i \neq i_0, j_0\) and \(j \neq i_0, j_0\) . If \(n > 2\) , there exist, for all distinct indices \(i_0\) and \(j\) , distinct indices \(i\) and \(j_0\) such that \(i \neq i_0, j_0 \neq j, j_0 \neq i_0\) ; hence \(u_{i_0j} = 0\) . This proves that an element of the centre of \(\mathfrak{g}\) is zero.
\item
  When \(\beta\) is alternating and \(n = 2m\) , \(\beta\) can be identified with the bilinear form on \(\mathbf{K}^{2m}\) with matrix \(\left( \begin{array}{cc}0 & I_m \\ -I_m & 0 \end{array} \right)\) with respect to the canonical basis (Algebra, Chapter IX, §5, Corollary 1 to Theorem 1). Under these conditions \(\mathfrak{g}\) is identified with the Lie algebra of matrices of the form \(U = \left( \begin{array}{cc}A & B\\ C & D \end{array} \right)\) where \(D = -^t A\) , \(B\) and \(C\) are symmetric ( \(A, B, C, D\) in \(\mathbf{M}_m(\mathbf{K})\) ) (§3, no. 4, Example 1). We use first the fact that \(U\) commutes with the matrix \(\left( \begin{array}{cc}X & 0 \\ 0 & -^t X \end{array} \right)\) , where \(X \in \mathbf{M}_m(\mathbf{K})\) . Then \(AX = X A\) , \(CX = -^t XC\) , \(XB = -B\) . \(^t X\) ; as these equalities must hold for all \(X\) , it follows that \(A\) is a scalar matrix \(\lambda I_m\) . We now use the fact that \(U\) commutes with the matrix \(\left( \begin{array}{cc}0 & Y \\ 0 & 0 \end{array} \right)\) , where
\end{enumerate}

\(Y\) is a symmetric matrix of \(\mathbf{M}_m(\mathrm{K})\) . Then \(\lambda Y = YC = CY = 0\) . This proves first that \(\lambda = 0\) . Moreover, for all \(X \in \mathbf{M}_m(\mathrm{K})\) , \(X + {}^t X\) is symmetric and hence \(XC = -{}^t XC\) . Using the equation \(CX = -{}^t XC\) obtained above, we see that \(C\) commutes with every element of \(\mathbf{M}_m(\mathrm{K})\) and hence that \(C\) is a scalar matrix, necessarily zero since \(YC = 0\) . It is similarly shown that \(B = 0\) .

For \(\beta\) symmetric and \(n = 2\) , \(\mathfrak{g}\) is of dimension 1 and hence commutative. For the other cases cf.~Exercises 25 and 26.

